\section{Introduction}
\label{sec:intro}

Multi-teacher on-policy distillation (MOPD) transfers domain specialists
into one student by scoring student-generated responses with their
assigned teachers \citep{ma2026mopdmultiteacheronpolicydistillation}.
The objective is to integrate their capabilities, but equal numbers
of prompts from each domain need not give those domains equal influence
on training. Open-MOPD identifies capability imbalance and studies
token-share weighting, changing distillation gaps, and stale
student-dependent rewards \citep{gao2026openmopddiagnosingfixingcapability}.
These findings motivate examining how the objective weights each
domain and which student parameters receive its strongest gradients.

Prior work reports sparse cumulative parameter changes in RL
\citep{mukherjee2025subnetworks} and OPD
\citep{yu2026denseupdates,shen2026opdgeometry}.
Such observations do not establish that instantaneous gradients are
sparse: optimizer state and accumulation over training can change
both the size and location of parameter movement. We make this
distinction explicit and ask how multiple teachers act on a common
student during MOPD. Gradients measured in that shared student are
more directly relevant to this question than endpoints of separately
trained students.

After introducing the MOPD learning objective in
Section~\ref{sec:setting}, we develop three studies and place each
experiment alongside its corresponding question:
\begin{enumerate}
\item \textbf{Token balancing (Section~\ref{sec:normalization}).}
We derive how global-token, domain-token, and domain-response averaging
weight domains and response lengths, then compare their effects on
capability balance. The analysis explains existing loss reductions
without introducing a new allocation algorithm.
\item \textbf{Gradient sparsity (Section~\ref{sec:subnetworks}).}
We measure gradient concentration in single-teacher and joint OPD,
compare the parameter subsets receiving each teacher's strongest
gradients in the same MOPD student, and relate their overlap to
teacher distribution differences. Optimizer steps and cumulative
changes provide complementary measurements.
\item \textbf{Supervision density (Section~\ref{sec:density}).}
We compare sampled-token PG and teacher top-$k$ training in both
settings, with full-vocabulary gradients on common prefixes as a
local reference. We test whether gradient sparsity differs materially
and report capability alongside it.
\end{enumerate}

Overlapping gradient coordinates can carry compatible or opposing
changes, and sampled-token feedback still depends on all logits through
the softmax normalizer. Neither low overlap nor limited vocabulary
coverage therefore guarantees sparse gradients or better integration.
\missing{Replace the prospective empirical contributions with observed
findings and their measured scope once the experiments are complete.}
